\scanpage{88}
% archive sheet number: 87
de l'icosaèdre combinatoire.

Si nous admettons que $G$ est \inserted{fini} d'ordre $120$, comme le stabilisateur de $s_0$ contient le groupe \inserted{$\unsure{Z|\mathbb{Z}}$} engendré par $v$ (cyclique d'ordre $5$) en plus de $\sigma_1$ -- donc le groupe \unsure{diédral} $\{1,\sigma_1\}.Z$ (produit $1/2$ direct), on voit que le nb des sommets est un diviseur de $120/10=12$. Notons que les six points déjà écrits $s_0,s_1,s_2,s_3,\unsure{s'_3|s_3^*},s''_3$ sont \uline{distincts}, car leurs coordonnées sont
\diagram{\((11)\quad\left\{\begin{gathered}\begin{tikzcd}[ampersand replacement=\&,row sep=small,column sep=small]
s_0=(0,0,0)\arrow[d,bend left=60,leftrightarrow]\\
s_1=(\unsure{1|0},0,0)\\
s_2=(0,1,0)\arrow[d,bend left=60,leftrightarrow]\\
s_3=(0,0,1)\\
s'_3=(-(1+\alpha),1+\alpha,1)=(\deleted{\unsure{+}}\alpha',-\alpha',1)\\
s''_{\unsure{2|3}}=(0,1+\alpha,1)=(0,-\alpha',1)
\end{tikzcd}\end{gathered}\right.\)}
\margin{$\alpha'=-(\alpha+1)=-\alpha^{-1}$ (``l'autre racine'' de $\alpha^2+\alpha-1=0$)}

\deleted{\unsure{Il}} Il est clair que $s_0$ à $s_3$ sont tous distincts, il faut \deleted{voir} les comparer à $s'_3,s''_3$ et \unsure{comparer} entre eux. Comme la dernière coordonnée de $s'_3,s''_3$ \deleted{sont} est $1$, ils sont $\ne s_0,s_1,s_2$ dont la dernière coordonnée est nulle. Ils \deleted{ne peuvent être}\inserted{sont} égaux \inserted{l'un ou l'autre à} $s_3$ \deleted{\unsure{que ssi|que si}} $1+\alpha=0$\inserted{(i.e. $\alpha=-1$)}, qui est \inserted{donc} aussi la condition néc et suff. pour que $s'_3=s''_3$. Or \deleted{\unsure{$1+\alpha=0$}} si $\alpha=-1$
